% Mathématiques 5e — cours V3
% Nombres rationnels et fractions

\section{Objectifs}

\begin{itemize}
\item Comprendre qu'une fraction est un nombre.
\item Comparer des fractions.
\item Produire des fractions égales.
\item Mettre deux fractions au même dénominateur.
\item Additionner et soustraire des fractions de dénominateurs quelconques.
\item Résoudre des problèmes mobilisant des additions et soustractions de fractions.
\item Relier fraction, quotient et différentes écritures d'un même nombre.
\end{itemize}

\section{1. Fraction et quotient}

Une fraction :
\[
\dfrac{a}{b}
\]
avec :
\[
b\neq0
\]
représente le quotient de $a$ par $b$.

Ainsi :
\[
\dfrac34
\]
est le nombre obtenu en divisant $3$ par $4$.

On peut écrire :
\[
\dfrac34=0{,}75.
\]

\begin{definition}
Un nombre rationnel est un nombre qui peut s'écrire comme quotient de deux entiers :
\[
\dfrac{a}{b},
\qquad b\neq0.
\]
\end{definition}

\section{2. Une fraction est un nombre}

Une fraction n'est pas seulement une « part de gâteau ».

Elle possède une position sur une droite graduée.

\coursefigure{/images/mathematiques/5eme/rationnels-droite.svg}{Droite graduée en quarts comportant plusieurs fractions entre zéro et deux.}{Une fraction peut être inférieure, égale ou supérieure à un ; elle représente toujours un nombre.}

Par exemple :
\[
\dfrac34<1.
\]

Mais :
\[
\dfrac54>1.
\]

Et :
\[
\dfrac84=2.
\]

\section{3. Fractions supérieures à un}

Une fraction dont le numérateur est supérieur au dénominateur peut être supérieure à $1$.

Par exemple :
\[
\dfrac{17}{5}.
\]

On peut effectuer la division euclidienne :
\[
17=3\times5+2.
\]

Donc :
\[
\dfrac{17}{5}
=
3+\dfrac25.
\]

Cette écriture permet de mieux situer le nombre :
\[
3<\dfrac{17}{5}<4.
\]

\section{4. Fractions égales}

Une même valeur peut posséder plusieurs écritures fractionnaires.

\begin{propriete}
Multiplier le numérateur et le dénominateur d'une fraction par un même nombre non nul ne change pas sa valeur :
\[
\dfrac ab=\dfrac{a\times k}{b\times k}.
\]
\end{propriete}

Par exemple :
\[
\dfrac23
=
\dfrac{2\times5}{3\times5}
=
\dfrac{10}{15}.
\]

De même :
\[
\dfrac47
=
\dfrac{8}{14}.
\]

\section{5. Simplifier une fraction}

On peut utiliser la propriété précédente dans l'autre sens.

Par exemple :
\[
\dfrac{12}{18}.
\]

Le numérateur et le dénominateur sont divisibles par $6$ :
\[
\dfrac{12}{18}
=
\dfrac{12\div6}{18\div6}
=
\dfrac23.
\]

En 5e, simplifier des fractions repose principalement sur des facteurs et multiples connus.

\section{6. Comparer deux fractions de même dénominateur}

Si deux fractions ont le même dénominateur, on compare leurs numérateurs.

Par exemple :
\[
\dfrac37<\dfrac57.
\]

Les deux nombres utilisent la même unité de fraction :
\[
\dfrac17.
\]

Cinq septièmes sont donc plus grands que trois septièmes.

\section{7. Comparer deux fractions de même numérateur}

Si deux fractions positives ont le même numérateur, la fraction au plus petit dénominateur est la plus grande.

Par exemple :
\[
\dfrac34>\dfrac37.
\]

En effet, partager l'unité en $4$ parts produit des parts plus grandes que la partager en $7$ parts.

\section{8. Comparer en se ramenant au même dénominateur}

Pour comparer :
\[
\dfrac34
\]
et :
\[
\dfrac56,
\]
on cherche un dénominateur commun.

Un multiple commun de $4$ et $6$ est :
\[
12.
\]

On écrit :
\[
\dfrac34=\dfrac9{12},
\]
et :
\[
\dfrac56=\dfrac{10}{12}.
\]

Comme :
\[
9<10,
\]
on obtient :
\[
\boxed{\dfrac34<\dfrac56}.
\]

\section{9. Comparer à un nombre simple}

On peut parfois éviter la mise au même dénominateur.

Par exemple :
\[
\dfrac78<1
\]
car :
\[
7<8.
\]

Et :
\[
\dfrac{11}{9}>1
\]
car :
\[
11>9.
\]

Comparer à :
\[
1
\]
ou à :
\[
\dfrac12
\]
peut fournir un repère efficace.

\section{10. Additionner des fractions de même dénominateur}

Lorsque le dénominateur est identique :
\[
\dfrac ab+\dfrac cb
=
\dfrac{a+c}{b}.
\]

Par exemple :
\[
\dfrac35+\dfrac45
=
\dfrac75.
\]

\begin{attention}
On n'additionne pas les dénominateurs :
\[
\dfrac35+\dfrac45\neq\dfrac79.
\]
\end{attention}

Le dénominateur représente la taille de l'unité fractionnaire ; il reste donc identique.

\section{11. Soustraire des fractions de même dénominateur}

De même :
\[
\dfrac ab-\dfrac cb
=
\dfrac{a-c}{b}.
\]

Par exemple :
\[
\dfrac{11}{8}-\dfrac38
=
\dfrac88
=
1.
\]

\section{12. Pourquoi faut-il un dénominateur commun ?}

Additionner :
\[
\dfrac23+\dfrac14
\]
directement n'a pas de sens car les parts ne sont pas de même taille.

Les tiers et les quarts doivent d'abord être exprimés dans une unité commune.

\coursefigure{/images/mathematiques/5eme/rationnels-denominateur.svg}{Bandes fractionnées montrant la mise au même dénominateur avant une addition.}{Pour additionner des fractions de dénominateurs différents, on transforme les écritures afin de disposer de parts de même taille.}

\section{13. Chercher un dénominateur commun}

Pour :
\[
\dfrac23+\dfrac14,
\]
on cherche un multiple commun de :
\[
3
\]
et :
\[
4.
\]

On peut choisir :
\[
12.
\]

Alors :
\[
\dfrac23
=
\dfrac8{12},
\]
et :
\[
\dfrac14
=
\dfrac3{12}.
\]

Donc :
\[
\dfrac23+\dfrac14
=
\dfrac8{12}+\dfrac3{12}
=
\dfrac{11}{12}.
\]

\section{14. Exemple développé : addition}

Calculons :
\[
\dfrac56+\dfrac7{15}.
\]

Un dénominateur commun à $6$ et $15$ est :
\[
30.
\]

On transforme :
\[
\dfrac56=\dfrac{25}{30},
\]
et :
\[
\dfrac7{15}=\dfrac{14}{30}.
\]

Donc :
\[
\dfrac56+\dfrac7{15}
=
\dfrac{25}{30}+\dfrac{14}{30}
=
\dfrac{39}{30}.
\]

On peut simplifier par $3$ :
\[
\dfrac{39}{30}
=
\dfrac{13}{10}.
\]

Donc :
\[
\boxed{\dfrac56+\dfrac7{15}=\dfrac{13}{10}}.
\]

\section{15. Exemple développé : soustraction}

Calculons :
\[
\dfrac79-\dfrac5{12}.
\]

Un dénominateur commun est :
\[
36.
\]

On écrit :
\[
\dfrac79=\dfrac{28}{36},
\]
et :
\[
\dfrac5{12}=\dfrac{15}{36}.
\]

Donc :
\[
\dfrac79-\dfrac5{12}
=
\dfrac{28}{36}-\dfrac{15}{36}
=
\dfrac{13}{36}.
\]

Ainsi :
\[
\boxed{\dfrac79-\dfrac5{12}=\dfrac{13}{36}}.
\]

\section{16. Choisir un dénominateur efficace}

Il n'est pas obligatoire de choisir le plus petit dénominateur commun, mais un choix adapté simplifie les calculs.

Pour :
\[
\dfrac34+\dfrac5{12},
\]
on pourrait choisir :
\[
48.
\]

Mais :
\[
12
\]
est déjà un multiple de $4$.

Il est donc plus simple d'écrire :
\[
\dfrac34=\dfrac9{12}.
\]

Puis :
\[
\dfrac9{12}+\dfrac5{12}
=
\dfrac{14}{12}
=
\dfrac76.
\]

\section{17. Résoudre un problème de partage}

Une randonnée comporte trois étapes.

La première représente :
\[
\dfrac38
\]
du parcours.

La deuxième représente :
\[
\dfrac14.
\]

La part parcourue après deux étapes vaut :
\[
\dfrac38+\dfrac14.
\]

On transforme :
\[
\dfrac14=\dfrac28.
\]

Donc :
\[
\dfrac38+\dfrac28
=
\dfrac58.
\]

Il reste :
\[
1-\dfrac58
=
\dfrac38.
\]

La troisième étape représente :
\[
\boxed{\dfrac38}
\]
du parcours.

\section{18. Résoudre un problème de durée}

Un élève consacre :
\[
\dfrac25
\]
d'une heure à un exercice puis :
\[
\dfrac14
\]
d'une heure à une révision.

Le temps total vaut :
\[
\dfrac25+\dfrac14.
\]

Un dénominateur commun est :
\[
20.
\]

Ainsi :
\[
\dfrac25=\dfrac8{20},
\qquad
\dfrac14=\dfrac5{20}.
\]

Donc :
\[
\dfrac8{20}+\dfrac5{20}
=
\dfrac{13}{20}.
\]

Le temps total représente :
\[
\boxed{\dfrac{13}{20}\text{ h}}.
\]

\section{19. Relier fraction et pourcentage}

Une fraction peut parfois être transformée facilement en pourcentage.

Par exemple :
\[
\dfrac14=25\%.
\]

De même :
\[
\dfrac12=50\%.
\]

Et :
\[
\dfrac35
=
\dfrac{60}{100}
=
60\%.
\]

Ces changements d'écriture permettent d'adapter la représentation au contexte.

\section{20. Relier fraction et écriture décimale}

Certaines fractions possèdent une écriture décimale finie.

Par exemple :
\[
\dfrac58=0{,}625.
\]

D'autres ne possèdent pas d'écriture décimale finie.

Par exemple :
\[
\dfrac13=0{,}333\ldots
\]

La fraction :
\[
\dfrac13
\]
reste alors une écriture exacte pratique.

\section{21. Méthode de comparaison}

\begin{methode}
Pour comparer deux fractions :
\begin{enumerate}
\item vérifier si elles ont déjà le même dénominateur ;
\item chercher si une comparaison à $0$, $1$ ou $\dfrac12$ suffit ;
\item sinon, les écrire avec un dénominateur commun ;
\item comparer les numérateurs obtenus.
\end{enumerate}
\end{methode}

\section{22. Méthode d'addition ou soustraction}

\begin{methode}
Pour calculer :
\[
\dfrac ab\pm\dfrac cd,
\]
on :
\begin{enumerate}
\item cherche un dénominateur commun ;
\item transforme chaque fraction en une fraction égale ;
\item additionne ou soustrait les numérateurs ;
\item conserve le dénominateur commun ;
\item simplifie si possible ;
\item contrôle la vraisemblance.
\end{enumerate}
\end{methode}

\section{23. Contrôler un résultat}

Supposons :
\[
\dfrac34+\dfrac23.
\]

Les deux fractions sont supérieures à :
\[
\dfrac12.
\]

Leur somme doit donc être supérieure à :
\[
1.
\]

Un résultat comme :
\[
\dfrac5{12}
\]
serait immédiatement impossible.

Le contrôle par ordre de grandeur fonctionne donc aussi avec des fractions.

\section{24. Erreurs fréquentes}

\begin{itemize}
\item Additionner les numérateurs et les dénominateurs.
\item Modifier uniquement le dénominateur sans modifier le numérateur.
\item Chercher systématiquement un dénominateur énorme.
\item Confondre une fraction avec deux nombres séparés.
\item Oublier qu'une fraction peut être supérieure à $1$.
\item Simplifier une somme en divisant certains termes seulement.
\item Donner une réponse sans vérifier son ordre de grandeur.
\end{itemize}

\section{À retenir}

\begin{aretenir}
Une fraction est un nombre :
\[
\boxed{\dfrac ab,\quad b\neq0}.
\]

Deux fractions peuvent représenter le même nombre :
\[
\boxed{\dfrac ab=\dfrac{ka}{kb}}.
\]

Pour comparer, additionner ou soustraire des fractions de dénominateurs différents, on peut les écrire avec un dénominateur commun.

Pour des fractions de même dénominateur :
\[
\boxed{
\dfrac ab+\dfrac cb=\dfrac{a+c}{b}
}
\]
et :
\[
\boxed{
\dfrac ab-\dfrac cb=\dfrac{a-c}{b}
}.
\]

Le sens du résultat doit toujours être contrôlé.
\end{aretenir}
